\section{A Physically Realizable Fourier Transform}
\label{sec:physical}

Here ``physically realizable'' has a precise ideal linear meaning.  A signed
state is represented by two nonnegative rails; local transport matrices have
nonnegative entries and preserve a declared mass or constant potential; and an
observable is recovered by rail difference.  A second realization uses convex
equilibrium junctions with an explicit quadratic energy.  These element laws
are sufficient to realize every normalized Bruun stage.

\subsection{Positive-cone lift}

For a real matrix $M\in\R^{m\times n}$, define $M^+=\max(M,0)$ and
$M^-=\max(-M,0)$ entrywise, and set
\begin{equation}
 C(M)=\begin{bmatrix}M^+&M^-\\M^-&M^+\end{bmatrix},
 \qquad
 \Pi_n=\begin{bmatrix}I_n&-I_n\end{bmatrix}.
 \label{eq:cone-lift}
\end{equation}

\begin{theorem}[Exact cone projection]
For every real matrix $M$,
\begin{equation}
 \Pi_m C(M)=M\Pi_n.
 \label{eq:cone-intertwining}
\end{equation}
Consequently, for any composable sequence $M_1,\ldots,M_s$,
\begin{equation}
 \Pi C(M_s)\cdots C(M_1)=M_s\cdots M_1\Pi.
 \label{eq:cone-composition}
\end{equation}
\end{theorem}
\begin{proof}
Since $M=M^+-M^-$,
\[
 \Pi_m C(M)=[M^+-M^-,\;M^--M^+]=M\Pi_n.
\]
Repeated substitution proves \eqref{eq:cone-composition}.
\end{proof}

When dimensions are clear, subscripts on $\Pi$ are suppressed.  The cone lift
maps $\R_+^{2n}$ to $\R_+^{2m}$ and stores additional common mode in
$\ker\Pi_n=\{(u,u):u\in\R^n\}$.  Equal
amounts may be removed from both rails without changing the signed observable.
Composition is exact after projection even though the lifted matrices need not
multiply as $C(AB)$.

Define rail canonicalization componentwise by
\begin{equation}
 K(p,n)=(p-r,n-r),\qquad r_i=\min(p_i,n_i),
 \label{eq:rail-canonicalization}
\end{equation}
and signed injection by $J(x)=(x^+,x^-)$.  Canonicalization gives a literal
fixed-point law for the nonnegative state.

\begin{theorem}[Canonical rail fixed points]
The map $K$ is idempotent, $\Pi K=\Pi$, and
\begin{equation}
 \operatorname{Fix}(K)=\{(p,n):p_i n_i=0\text{ for every }i\}.
 \label{eq:canonical-fixed-set}
\end{equation}
For every real $M$ and signed $x$,
\begin{equation}
 KC(M)J(x)=J(Mx).
 \label{eq:canonical-stage}
\end{equation}
Consequently, inserting $K$ after every lifted Fourier stage yields the
canonical injection of the exact signed stage product.
\end{theorem}
\begin{proof}
After subtracting $r_i$, at least one rail in each pair is zero; a second
application subtracts zero, proving idempotence and
\eqref{eq:canonical-fixed-set}.  Equal subtraction preserves rail difference.
Equation \eqref{eq:cone-intertwining} shows that $C(M)J(x)$ has signed
projection $Mx$; canonicalization selects the unique disjoint-support
representative of that projection, namely $J(Mx)$.  Induction proves the stage
product statement.
\end{proof}

\subsection{Conservative mass gauge}

Consider a factorization $x_{s+1}=M_sx_s$.  Fix $g_S=\bm{1}$ and recurse backward
by
\begin{equation}
 g_s=|M_s|^{\trans}g_{s+1},
 \qquad
 A_s=\diag(g_{s+1})M_s\diag(g_s)^{-1}.
 \label{eq:mass-gauge}
\end{equation}

\begin{theorem}[Column-stochastic cone stages]
If every component of $g_s$ is positive, then $C(A_s)$ is nonnegative and
column-stochastic.  The diagonal gauges telescope, and terminal signed
projection with the declared boundary calibration recovers
$M_{S-1}\cdots M_0$ exactly.
\end{theorem}
\begin{proof}
For column $j$,
\[
 \sum_i|(A_s)_{ij}|=
 {\sum_i(g_{s+1})_i|(M_s)_{ij}|\over(g_s)_j}=1.
\]
Each column of $C(A_s)$ therefore sums to one.  With
$D_s=\diag(g_s)$, multiplying the gauged stages gives
\begin{equation}
 A_{S-1}\cdots A_0=D_S(M_{S-1}\cdots M_0)D_0^{-1}.
 \label{eq:mass-boundary-product}
\end{equation}
Injecting $J(D_0x)$ and decoding by $D_S^{-1}\Pi$ therefore recovers the
original signed product.  Here $g_S=\bm{1}$, so $D_S=I$.
\end{proof}

This construction preserves total two-rail mass at every transport stage.
Annihilation of equal rail amounts is a separately ledgered conversion into a
null observable; it cannot change any Fourier coordinate.

\subsection{Kinematic and equilibrium gauge}

For a nonnegative cone stage $C_s$ with no zero row, choose $d_0>0$ and define
componentwise
\begin{equation}
 d_{s+1}=\bigl(C_s(1/d_s)\bigr)^{-1},
 \qquad
 B_s=\diag(d_{s+1})C_s\diag(d_s)^{-1}.
 \label{eq:kinematic-gauge}
\end{equation}
Every normalized Bruun cell has a nonzero entry in each row, so direct sums and
coordinate permutations of its cone lifts satisfy the hypothesis.

\begin{proposition}[Convex junction form]
$B_s\bm{1}=\bm{1}$.  Hence each output coordinate is a convex combination of the
input coordinates and can be decomposed into binary interpolation junctions.
\end{proposition}
\begin{proof}
Substitution of \eqref{eq:kinematic-gauge} gives
$B_s\bm{1}=\diag(d_{s+1})C_s(1/d_s)=\bm{1}$.
\end{proof}

Writing $\Delta_s=\diag(d_s)$ gives the boundary-calibrated product
\begin{equation}
 B_{S-1}\cdots B_0
 =\Delta_SC_{S-1}\cdots C_0\Delta_0^{-1}.
 \label{eq:kinematic-boundary-product}
\end{equation}
Thus a cone input is encoded by $\Delta_0$, and the output is decoded by
$\Delta_S^{-1}$ before applying $\Pi$.

\begin{theorem}[Exact clamped equilibrium]
Let $B\ge0$ be row-stochastic and hold $x$ fixed.  The energy
\begin{equation}
 E_x(y)={1\over2}\sum_{j,i}B_{ji}(y_j-x_i)^2
 \label{eq:equilibrium-energy}
\end{equation}
has the unique equilibrium
\begin{equation}
 y=Bx.
 \label{eq:equilibrium-map}
\end{equation}
\end{theorem}
\begin{proof}
Row stochasticity gives
$\partial E_x/\partial y_j=y_j-\sum_iB_{ji}x_i$.  The Hessian with respect to
$y$ is $I$, so the stationary point \eqref{eq:equilibrium-map} is unique.
\end{proof}

The theorem applies to an isolated stage with prescribed input, or once to the
row-stochastic product of all kinematic stages.  The latter makes the complete
Fourier boundary response an exact static equilibrium.  The sparse exact
realization is instead the ideal kinematic-constraint or reciprocal-force
cascade.  A reciprocal linkage obeying
$q_{\rm in}=A^{\trans}q_{\rm out}$ also transmits generalized force as
$f_{\rm out}=Af_{\rm in}$ by equality of virtual work.  Thus the transpose
kinematic constraint and forward Fourier force map are the same reciprocal
linear element.

\subsection{Field-phase DIP form}

Before conjugate folding, the normalized complex butterfly is
\begin{equation}
 {1\over\sqrt2}\begin{bmatrix}1&w\\1&-w\end{bmatrix},
 \qquad |w|=1.
 \label{eq:lossless-butterfly}
\end{equation}
It is a fixed phase delay followed by a lossless balanced coupler, since the
matrix in \eqref{eq:lossless-butterfly} is unitary.  In the real two-plane
description, the phase delay is $R_\theta$ and addition is a linear junction.
Conjugation required by a folded half-spectrum is a separate real reflection.
Thus the unfurled DIP lattice has an exact passive scattering realization of
$N^{-1/2}\DFT_N$; a declared boundary gain $\sqrt N$ recovers the unnormalized
Fourier convention.  The folded construction adds the declared reflection
element.  The cone and field constructions are two ideal physical categories
for the same Fourier observable.

\begin{corollary}[Normalized Bruun realization]
Apply the cone lift and either diagonal gauge to any exact stage factorization
of $B_N$.  With the boundary encodings in
\eqref{eq:mass-boundary-product} or
\eqref{eq:kinematic-boundary-product}, rail-difference projection equals
$B_Nx$, and standard output permutation gives $\DFT_N^{\R}x$ exactly.
\end{corollary}
